\section{Importazione dati nei Database}
\label{sec:database_ingestion}

Una volta completata l'estrazione e la strutturazione dei documenti legali in formato JSON, si apre la fase di persistenza e indicizzazione dei dati all'interno dei databse. Nel contesto di questo progetto di tirocinio, l'architettura prevede il caricamento dei dati estratti all'interno di due distinte tipologie di database: Neo4j, un database a grafo, e Qdrant, un database vettoriale.

L'uso di due database consente di sperimentare l'uso di entrambi i paradigmi di memorizzazione sullo stesso dominio normativo, al fine di valutarne i rispettivi punti di forza, analizzare le performance e far emergere differenze qualitative e quantitative nelle fasi di interrogazione e navigazione dei testi legali.

Per garantire la coesistenza flessibile di queste due tecnologie e favorire un ambiente di sviluppo pulito e modulare, lo sviluppo del software ha visto l'applicazione combinata di due noti design pattern software: lo \textit{Strategy Pattern} e il \textit{Repository Pattern}.

\subsection{I Modelli di Dominio}
Prima di analizzare le strategie di persistenza, è fondamentale descrivere come i dati estratti vengano rappresentati in memoria. All'interno della cartella \texttt{models}, il sistema definisce le entità fondamentali del dominio normativo utilizzando il paradigma ad oggetti. Queste classi operano essenzialmente come \textit{Models}: strutture dati pure, prive di logica di persistenza diretta o dipendenze verso i DBMS, il cui unico scopo è incapsulare i dati e preservare i vincoli di integrità del dominio.

L'architettura dei modelli rispecchia fedelmente la scomposizione gerarchica del testo legale:
\begin{itemize}
    \item \texttt{Entity}: Rappresenta l'ente pubblico emittente (ad esempio, un Ministero o un Comune), memorizzando i metadati identificativi dell'istituzione.
    \item \texttt{LegalAct}: Modella l'atto normativo nella sua interezza. Al suo interno, oltre ai metadati generali del documento, espone il metodo di fabbrica \texttt{from\_json\_file(file\_path)}, delegato al parsing del file JSON e alla deserializzazione ricorsiva dei suoi elementi strutturali.
    \item \texttt{StructuralElement}: Una classe base polimorfica specializzata nelle sottoclassi \texttt{Title}, \texttt{Chapter} e \texttt{Article}. Permette di mappare i livelli intermedi di aggregazione del testo, mantenendo traccia del numero d'ordine e del titolo della sezione.
    \item \texttt{Paragraph}: Rappresenta l'unità atomica informativa e testuale, ovvero il singolo comma o blocco di testo non ulteriormente scindibile, sul quale verranno poi calcolati gli embedding.
\end{itemize}

\subsection{Disaccoppiamento della Persistenza: Lo Strategy Pattern}
L'esigenza di effettuare una doppia ingestion prima su Neo4j e poi su Qdrant senza dover duplicare la logica di orchestrazione o modificare il flusso principale del programma ha reso naturale l'adozione del design pattern \textit{Strategy}. Questo pattern comportamentale definisce una famiglia di algoritmi, ne incapsula ciascuno all'interno di una classe e li rende tra loro intercambiabili, consentendo all'algoritmo di variare indipendentemente dai client che lo utilizzano.

Nel sistema sviluppato, l'interfaccia astratta di riferimento è rappresentata dalla classe \texttt{DBImporterStrategy}. Come mostrato nel listato del codice sorgente, essa definisce tre metodi fondamentali che ogni classe concreta deve implementare:
\begin{itemize}
    \item \texttt{connect()}: per la gestione e la validazione della connessione con l'istanza specifica del database;
    \item \texttt{import\_legal\_act(entity, legal\_act)}: per l'esecuzione materiale della logica di Ingestion dell'atto normativo e dei metadati dell'ente emittente;
    \item \texttt{close()}: per la chiusura sicura delle sessioni.
\end{itemize}

Nello script principale \texttt{main.py}, il ciclo di esecuzione batch istanzia dinamicamente l'oggetto concreto (\texttt{Neo4jImporter} o \texttt{QdrantImporter}) a seconda della configurazione desiderata. La funzione \texttt{process\_documents} riceve l'istanza polimorfica tramite la \textit{Dependency Injection} ed esegue il parsing e l'ingestion dei file JSON senza curarsi di quale sia il database di destinazione. Questa architettura estendibile riduce a zero l'accoppiamento e semplifica future implementazioni con ulteriori database.


\subsection{Modularità dell'Accesso ai Dati: Il Repository Pattern}
Se lo \textit{Strategy Pattern} isola i database nella loro interezza rispetto all'orchestrazione principale, il \textit{Repository Pattern} interviene specificamente all'interno dell'architettura di persistenza dei database per mediare tra lo strato dei modelli di dominio e lo strato di accesso ai dati fisici. 

L'obiettivo principale del Repository Pattern è centralizzare la logica di costruzione ed esecuzione delle query, offrendo al client un'interfaccia astratta simile a quella di una collezione di oggetti in memoria. Nel contesto di Neo4j, un atto legale non può essere salvato tramite un'unica operazione atomica standard, a causa delle profonde interconnessioni tra i nodi della gerarchia. Anziché concentrare la scrittura di stringhe Cypher complesse all'interno della classe \texttt{Neo4jImporter}, la responsabilità è stata distribuita su una serie di repository specializzati, ognuno deputato alla gestione del ciclo di vita di un singolo modello di dominio:
\begin{itemize}
    \item \texttt{EntityRepository}: Gestisce l'inserimento e l'indicizzazione dei nodi relativi agli enti istituzionali.
    \item \texttt{LegalActRepository}: Inserisce il nodo radice dell'atto e lo ancora all'ente emittente.
    \item \texttt{StructuralElementRepository}: Incapsula la logica Cypher necessaria a creare i nodi intermedi del grafo e a tessere gli archi gerarchici verso i rispettivi nodi padri.
    \item \texttt{ParagraphRepository}: Si occupa esclusivamente della persistenza dei nodi foglia contenenti il testo dei commi.
\end{itemize}

In questo modo, l'importer concreto delega l'interazione fisica con il database ai singoli repository, i quali ricevono il contesto transazionale attivo e l'oggetto di dominio da persistere, nascondendo completamente i dettagli del linguaggio Cypher al resto dell'applicazione.

\subsection{Dettagli Implementativi dell'Ingestione su Neo4j}
L'implementazione racchiusa in \texttt{Neo4jImporter} si concentra sulla fedele ricostruzione topologica e gerarchica del documento all'interno del grafo. Per garantire l'integrità transazionale e prevenire stati inconsistenti dovuti a fallimenti parziali della rete o del backend, il caricamento di ogni singolo documento viene eseguito all'interno di un'unica transazione ACID esplicita.

All'interno della funzione transazionale di lavoro, viene orchestrato un flusso sequenziale:
\begin{enumerate}
    \item Viene invocato l'\texttt{EntityRepository} per assicurare la presenza del nodo dell'ente nel grafo.
    \item Il \texttt{LegalActRepository} crea il nodo radice del documento normativo, collegandolo all'ente tramite una relazione orientata.
    \item L'importer itera sulla lista lineare degli elementi strutturali pre-ordinati estratti dal modello (\texttt{legal\_act.elements}). In questa fase si applica il polimorfismo a runtime: tramite costrutti di introspezione (\texttt{isinstance}), il sistema smista l'elemento al repository corretto. Se l'oggetto è un'istanza di \texttt{StructuralElement} viene invocato lo \texttt{struct\_repo}, mentre se si tratta di un \texttt{Paragraph} interviene il \texttt{par\_repo}. Ciascuna chiamata riceve l'identificativo del nodo padre (\texttt{parent\_id}) e la tipologia di arco (\texttt{relation\_type}), consentendo di generare relazioni stabili come \texttt{HAS\_CHAPTER} o \texttt{HAS\_PARAGRAPH}.
\end{enumerate}

Per qunato riguarda l'ingestione delle relazioni trasversali (es. citazioni ad altri articoli), queste relazioni vengono lette da file JSON esterni dedicati. 
Il metodo \texttt{import\_relationships} esegue una seconda passata sul grafo in cui, per ogni link estratto, esegue una query Cypher che sfruttando l'indicizzazione sul \texttt{custom\_id}, individua istantaneamente i due nodi (sorgente e destinazione) e traccia l'arco semantico specifico.

\subsection{Dettagli Implementativi dell'Ingestione su Qdrant}
La strategia implementata nella classe \texttt{QdrantImporter} risponde a logiche radicalmente differenti, dettate dai vincoli operativi dei database vettoriali e dalle necessità algoritmiche della ricerca semantica (\textit{Semantic Search}). Qdrant non supporta nativamente strutture gerarchiche ramificate, ma organizza le informazioni in collezioni piatte di vettori denominati \textit{Points}.

Il fulcro del processo di ingestion in Qdrant risiede nell'operazione di appiattimento della struttura ad albero dell'atto normativo. L'unità fondamentale di indicizzazione viene identificata nel modello \texttt{Paragraph}, in quanto porzione di testo ottimale per bilanciare la granularità informativa e la focalizzazione semantica del modello di embedding. Tuttavia, memorizzare un comma in modo isolato causerebbe una perdita critica di contesto.
Per risolvere questo problema, l'algoritmo esegue una procedura di \textit{tree-climbing} in memoria per ciascun paragrafo: partendo dall'elemento foglia, risale ricorsivamente la catena dei nodi padri, durante la risalita, estrae il numero e il titolo di ogni livello gerarchico attraversato (Articolo, Capitolo, Titolo) e tutte queste informazioni vengono aggregate e inserite all'interno del dizionario di metadati del punto vettoriale, denominato \textit{Payload}. Il payload risultante conterrà quindi l'identificativo univoco, il testo del comma e l'intero cammino strutturale di provenienza.